\documentclass[11pt]{ekblite}
\newcommand\mtop{1in}
\newcommand\mbottom{1in}
\newcommand\mleft{1.2in}
\newcommand\mright{1.2in}
\usepackage[top=\mtop, bottom=\mbottom, left=\mleft, right=\mright]{geometry}
\pagenumbering{arabic}

\newcommand\aaron[1]{\textcolor{red}{#1}}
\hypersetup{ 	
pdfsubject = {blog},
pdftitle = {thinking},
pdfauthor = {Aaron Huntley}
}

\begin{document}

\title{Thinking}
\maketitle

%intro

Can you ever have a unique thought? 
Maybe if you have a thought which is long and complex enough then yes, but what makes this thought unique in an interesting way?
For example you can think of a really long string of numbers and its increasingly unlikely that no one has thought of that exact string of numbers.
However, that thought doesn't seem interesting.
So what makes a thought interesting?

How can we separate our thoughts into a single thought particle?
At least how my mind works is a stream of thoughts never ending each one influencing the next.
In this way are all thoughts unique? 
Each mind has its own unique interpretation of thoughts and though we can transfer the thought information through communication is a thought ever perfectly replicated?
How could it possibly be perfectly replicated without exactly the same hardware.
Then maybe we are being too strict.
We should consider thoughts up to some looser form of equality.
If you think of a red car and I think of a red car, we must both be thinking of the same thing?
Even while you read this now, it is an attempt for me to make you think what I am thinking of, and yet we will be thinking of these things differently.
I like this, because now we get two thoughts for the price of one.

Each person we speak to can offer a new perspective on a thought. 
Can we get a complete picture of a single thought if we collect all the perspectives.
Perhaps we can define humanities collective thought this way.
Then humanities collective still a finite number of thoughts, but is what were trying to think infinite?
Is a thought finite?
You can argue the space of all thoughts is at least countably infinite.
Since if you think of $1$ then you can think of $2$, and by induction all $n\in\mathbb{N}$.
However, when I think of a red car is that one thought.
Is the instance of the electrons firing in my brain enough to define that thought?

I believe I am only capable of thinking of a finite number of things in the span of my life.
If this is the case then what we think about is scarce.
So if we can choose what to think about then what so we spend our thoughts on?
Do you want to have the most unique thoughts?
Do you like thinking about the same things?
Do you only want to think about things that bring you happiness?

Can you think of anything you want?
If I tell you to do that what does your mind go to?
Mine goes to something concrete, like a giraffe with a really long neck, just to prove to my self I can think of anything
However, I don't think I can.
When I ask my brain to think of anything, it tries to fill the gap but I feel there are things it will never think.
So when I think I like to wonder in the dark forest of thought and observe whats there.
I like to do this in conversation with others, since they bring their own lantern to the forest and shine it on the leaves, trees and wildlife which I don't see.
Each time you wonder the trail of a thought, you may be familiar with whats there but you could spot something new or interesting, something which makes you smile.
As in life.

thoughs vs potential thogubts vs possible thoughts. finiteness of these.

\section{part - what to think about}

Once you come to the realisation that thoughts are finite, it is compelling to try to change your thoughts to try to think about things we 'want' to.
You will be quick to realise that this is not as easy as it first seems.
We easily get distracted by what we want for dinner, the scratch on your thigh and oh the doorbell just rang.
What we think about is a consequence of both our environment and what we feel.
In this case, at the very least can we observe these thoughts?
Common practices of meditation try to eliminate the environmental factor to observe both state of mind and body.
By focusing on something like the breath you can watch your own mind working as an organ in your body, as surely as your heart beats the mind thinks.
Aside from studies showing that mediation actually does change your brain structure \aaron{cite}, this process \emph{feels} as if you discover what it means to control your thoughts.
In fact it feels more like you can't control them but just observe them, perhaps guide them.
Now it certainly feels that we can make choices to do things, I agree with this, and those choices inturn affect our thoughts.
However this is a more indirect way of 'controlling' thoughts.

In meditation we can also observe feelings, to us feelings and thoughts may be both distinct and intertwined. 
A feeling may come with something more physical or emotional but thoughts may seem more objective and grounded.
However, sometimes it may be hard to separate the two, feelings can affect thoughts and thought can affect feelings.
But just as with thoughts we don't have much control over feelings.
We can merely observe them.
Again, via something such as meditation, observing them has some benefit.
However, it does not give us complete control.

Okay so what do we do?
choose morals/hobbies/joy eg flowers
do things based on those
obvserve ourselves doiung those things.





\end{document}
